% ============================================================
% appendix.tex -- Appendices A, B, C (the main file places \appendix before this).
% Numbers come from raw_materials/experimental_log.md. Appendix A additionally
% states commit ids and preregistered decision rules taken from the frozen
% preregistration texts (branch research/patent-citation-hole-20261118).
% Revised in round 1 (REVISION_PLAN_round1.md, sections A--C).
% ============================================================

\section{Preregistration and Amendment Log}
\label{app:prereg}

Study~2 rests on four preregistration plans. Each is a self-timestamped git commit in PR~\#5 of the project repository; none was registered with a third-party registry such as OSF, and the commit timestamps are set by the authors' own tooling. Preregistration is meant to separate confirmatory from exploratory analysis by fixing the plan before outcomes are observed \citep{nosek2018preregistration}. By that standard, Plan~1 was fixed before any patent data were downloaded. Plans 2--4 were written after earlier results were known, so the questions they ask were data-informed; only their definitions, null models, decision rules, thresholds and stop rules were fixed before the analyses they govern were computed. Changes after freezing are recorded as dated amendments that do not overwrite the frozen text. Table~\ref{tab:app_prereg} summarises the four plans, including the commit ids and rule thresholds taken from the frozen texts; Table~\ref{tab:app_amend} lists every amendment and recorded deviation.

\begin{table}[tbp]
\centering
\caption{Preregistration log for Study~2. ``Frozen before'' states what had not yet been computed when the plan was committed, and ``Known when written'' what had. Commit ids and decision rules are quoted from the frozen texts; outcomes are DIRECT results (E3).}
\label{tab:app_prereg}
\footnotesize
\setlength{\tabcolsep}{4pt}
\begin{tabular}{p{2.0cm}p{2.5cm}p{3.7cm}p{4.3cm}p{2.3cm}}
\toprule
Plan (commit) & Frozen before / known when written & Hypotheses and content & Decision rule & Outcome \tabularnewline
\midrule
\raggedright 1.~Aggregate test (\texttt{a8f3289}) &
\raggedright Frozen before any patent data were downloaded or inspected. Known: Study~1 only &
\raggedright H1a: OE(AI,DEC) relative to reference CPC main groups in G06/G16/H04 with $\geq$2,000 patents; H1b: ratio to OE(AI,BIZ-OTHER); H2: trend by citing year; H3 (exploratory): assignees as brokers of a domain-level hole; six sensitivity checks; synthetic calibration &
\raggedright H1 supported only if OE is in the lowest quartile on both sides (H1a) and the ratio's upper 95\% CI bound is $<1$ (H1b); labels \texttt{HOLE\_\allowbreak SUPPORTED}, \texttt{PARTIAL\_\allowbreak NOT\_\allowbreak CLAIMED}, \texttt{HOLE\_\allowbreak NOT\_\allowbreak SUPPORTED}; $\geq$4 of 6 sensitivity checks disagree $\rightarrow$ \texttt{FRAGILE}; H2 labels \texttt{CLOSING}, \texttt{PERSISTENT}, \texttt{NARROWING\_\allowbreak BUT\_\allowbreak OPEN} &
\raggedright \texttt{HOLE\_\allowbreak NOT\_\allowbreak SUPPORTED}; 0/6 disagree; H2 \texttt{PERSISTENT} \tabularnewline
\addlinespace
\raggedright 2.~V9 battery on the CPC matrix (\texttt{a6266f7}) &
\raggedright Frozen before any cell-level count. Known: the aggregate verdict &
\raggedright Unchanged V9 engine on a $6\times7$ matrix fixed from the CPC scheme; Part~A co-classification, Part~B applicant citations &
\raggedright Full core: q $\leq$ .05 in main, exact and document-preserving tests in the same direction, plus raw; convergent depletion: full-core depleted in both A and B; a depleted cell may not be called a structural hole &
\raggedright Preregistered core 0 by construction (B1); three-test cells exploratory \tabularnewline
\addlinespace
\raggedright 3.~Full replication suite (\texttt{d641256}) &
\raggedright Frozen before any module P1--P10 was run. Known: the aggregate verdict, the Plan~2 full-window cell results and the M06 failure (B1) &
\raggedright P1 with the patent-adapted null M06-P; P2 seeds, null false positives, power, segmentation; P3 time split; P4 label face validity; P5 firm network; P6 removal; P7 hyperedges; P8 aliases; P9 Gap Validity; P10 Limitation Scout &
\raggedright Three-test agreement for P3 and P9; replicated depleted cell must pass in both halves; P4 pass if lift $\geq$1.5; P5 firm hole if O/E $<1$ and below both MIXED O/Es; a candidate is refuted if AI and DEC firms stay connected after its removal &
\raggedright See Appendix~\ref{app:sens}; P5 criterion met (later \texttt{PARTIAL} under the Plan~4 null; not a confirmed hole); 10/10 candidates refuted; P10 \texttt{NOT\_\allowbreak APPLICABLE} \tabularnewline
\addlinespace
\raggedright 4.~Portfolio-preserving firm null (\texttt{b59fc71}) &
\raggedright Frozen before any null replicate was computed. Known: the P5 result, which motivated the plan &
\raggedright T1: OE(AI\_FIRM$\sim$DEC\_FIRM); T2: MIXED median constraint; owner sets permuted within strata (primary; S-tech). The null preserves portfolio counts, not firm degree &
\raggedright \texttt{ORGANISATIONAL\_\allowbreak HOLE\_\allowbreak BEYOND\_\allowbreak PORTFOLIO} if T1 and T2 both have one-sided $p<.05$; \texttt{PARTIAL} if exactly one; \texttt{PORTFOLIO\_\allowbreak ARTEFACT} if neither; S-tech disagreement $\rightarrow$ \texttt{TECH\_\allowbreak MIX\_\allowbreak SENSITIVE} &
\raggedright \texttt{PARTIAL} in both stratifications (T1 met, T2 not) \tabularnewline
\bottomrule
\end{tabular}
\end{table}

\begin{table}[tbp]
\centering
\caption{Amendments and recorded deviations. Amendments are appended to the frozen texts, which remain unchanged.}
\label{tab:app_amend}
\footnotesize
\setlength{\tabcolsep}{4pt}
\begin{tabular}{p{1.3cm}p{1.2cm}p{2.6cm}p{6.4cm}p{3.1cm}}
\toprule
Item & Plan & When & Change & Effect on design \tabularnewline
\midrule
\raggedright A1 & \raggedright 1 & \raggedright 2026-11-18, before any citation data or result was computed (\texttt{a158e70}) & \raggedright The CPC file labels inventive codes ``inventional''; the code accepts both spellings so that sensitivity check s4 keeps its meaning. CPC rows are loaded in chunks. & \raggedright None \tabularnewline
\addlinespace
\raggedright A2 & \raggedright 1 & \raggedright 2026-11-18, after the run (data note) & \raggedright Before 2013, non-examiner citations are labelled ``cited by other'', so the applicant-only specification effectively covers citing years 2013--2025; s1 (all categories) agrees with the primary verdict. H3 was not run. H3 asked whether assignees broker a \emph{domain-level} hole, so it was conditional on H1; with H1 unsupported there was no such hole to broker. The later firm analysis asks a different question---whether specialist AI and DEC firms are separated, and whether mixed-portfolio firms span that separation, irrespective of the domain-level verdict. It was preregistered separately in Plan~3 (P5) as the patent analogue of the Study~1 author-network line, and followed up in Plan~4. & \raggedright None to H1; H3 moot; firm-level question re-posed in Plans 3--4 \tabularnewline
\addlinespace
\raggedright B1 & \raggedright 2 & \raggedright 2026-11-18, after the run & \raggedright M06 accepted 0 of 21,843,600 proposed swaps because each document is a complete row $\times$ column block; the preregistered core is 0 by construction and uninformative. Three-test cells are reported as exploratory. The replacement M06-P was specified only after this failure and after the Plan~2 three-test cell results were known; it was then frozen in Plan~3 (module P1, \texttt{d641256}) before it was computed. The Plan~2 result file records only the degenerate core; all M06-P counts come from the Plan~3 P1 output. & \raggedright Documented, data-informed deviation from Plan~2 \tabularnewline
\addlinespace
\raggedright Seeding & \raggedright 4 & \raggedright Before any null replicate finished (\texttt{1eb7319}) & \raggedright The first run was stopped by host memory pressure after reproducing P5 exactly. The run was made resumable; each replicate $r$ of scheme $s$ uses its own seed sequence [20261119, $s$, $r$] instead of one stream. & \raggedright Statistic, strata and rule unchanged \tabularnewline
\addlinespace
\raggedright P6, P10 & \raggedright 3 & \raggedright Recorded with results & \raggedright The P6 cross-weight-retained metric is identically 1 by construction; P10 Limitation Scout is \texttt{NOT\_\allowbreak APPLICABLE} (patents have no limitation section). & \raggedright Reported as uninformative \tabularnewline
\bottomrule
\end{tabular}
\end{table}

Study~1 used frozen V9 seeds: 20260918 for the stub Monte Carlo and 20260919 for the document-preserving swap (E1.1). Its data-validity re-analysis (2026-11-18; E1.3--E1.5) was not preregistered. Its verdict was recorded as follows: directional under-integration DEFENSIBLE; ``16 statistically robust depleted pairs'' NOT DEFENSIBLE as measured; Burt structural hole STOP; research gap NOT SUPPORTED. We qualify the first item in this paper. It rests on the attenuation argument, which holds only for non-differential misclassification, whereas same-term self-pairing and model-inferred labels are non-random errors. The direction survives every check we could run, but it is descriptive and unconfirmed (Section~\ref{sec:discussion_measurement}).

% ------------------------------------------------------------
\section{Sensitivity Tables}
\label{app:sens}

\subsection{Study 1}

Table~\ref{tab:app_v9cells} lists the 16 depleted cells of the frozen V9 analysis, and Table~\ref{tab:app_s1robust} the robustness checks of the data-validity re-analysis.

\begin{table}[tbp]
\centering
\caption{The 16 V9 depleted cells with observed counts (DIRECT; \texttt{V9\_ROBUSTNESS\_MASTER\_TABLE.csv}). These are depleted cells, not gaps, and they are withdrawn as anchors of the model (Section~\ref{sec:mlp_model}). Nine of the 16 lie in the \texttt{PREDICTION\_\allowbreak RISK\_\allowbreak ESTIMATION} row or the \texttt{PREDICTION\_\allowbreak DIAGNOSIS} column, where same-term self-pairing in the prediction $\times$ prediction cell can deplete cells mechanically under fixed margins.}
\label{tab:app_v9cells}
\footnotesize
\begin{tabular}{llr}
\toprule
AI family & Decision family & $O$ \\
\midrule
\texttt{COORDINATION\_EXECUTION} & \texttt{HUMAN\_REVIEW\_ACCOUNTABILITY} & 10 \\
\texttt{COORDINATION\_EXECUTION} & \texttt{RISK\_ASSESSMENT} & 3 \\
\texttt{EVALUATION\_ASSESSMENT} & \texttt{ALLOCATION\_PLANNING} & 12 \\
\texttt{EVALUATION\_ASSESSMENT} & \texttt{PREDICTION\_DIAGNOSIS} & 9 \\
\texttt{EXPLANATION\_ASSURANCE} & \texttt{ALLOCATION\_PLANNING} & 7 \\
\texttt{EXPLANATION\_ASSURANCE} & \texttt{EVALUATION\_SELECTION} & 7 \\
\texttt{EXPLANATION\_ASSURANCE} & \texttt{PREDICTION\_DIAGNOSIS} & 15 \\
\texttt{GENERATION\_SYNTHESIS} & \texttt{PREDICTION\_DIAGNOSIS} & 7 \\
\texttt{HUMAN\_AUGMENTATION} & \texttt{EVALUATION\_SELECTION} & 2 \\
\texttt{HUMAN\_AUGMENTATION} & \texttt{PREDICTION\_DIAGNOSIS} & 3 \\
\texttt{PREDICTION\_RISK\_ESTIMATION} & \texttt{EVALUATION\_SELECTION} & 8 \\
\texttt{PREDICTION\_RISK\_ESTIMATION} & \texttt{GENERAL\_DECISION\_PROCESS} & 8 \\
\texttt{PREDICTION\_RISK\_ESTIMATION} & \texttt{GOVERNANCE\_OVERSIGHT} & 26 \\
\texttt{PREDICTION\_RISK\_ESTIMATION} & \texttt{HUMAN\_REVIEW\_ACCOUNTABILITY} & 6 \\
\texttt{PREDICTION\_RISK\_ESTIMATION} & \texttt{INFORMATION\_ACQUISITION} & 21 \\
\texttt{RECOMMENDATION\_OPTIMIZATION} & \texttt{HUMAN\_REVIEW\_ACCOUNTABILITY} & 8 \\
\bottomrule
\end{tabular}
\end{table}

\begin{table}[tbp]
\centering
\caption{Study~1 robustness checks (DIRECT; data-validity re-analysis, 2026-11-18; E1.3). Core cells are defined by BH-FDR within each test \citep{benjamini1995controlling} combined with agreement across tests, so the core has no single defined false discovery rate. The checks test the stability of the pattern; none of them addresses non-random coding error.}
\label{tab:app_s1robust}
\footnotesize
\begin{tabular}{p{3.6cm}p{11.0cm}}
\toprule
Check & Result \tabularnewline
\midrule
\raggedright Re-implementation & \raggedright Exact and raw tests reproduce the official q values ($|\Delta q|\leq 6\times10^{-12}$); Jaccard overlap of the 32-cell core .94 \tabularnewline
\raggedright 20 Monte Carlo seeds & \raggedright 31 cells in the core under every seed, including all 16 depleted cells; 2 boundary cells, exactly the one-in / one-out pair of the human-corrected sidecar \tabularnewline
\raggedright Exclusion bias & \raggedright Depleted core cells pass the strict filter more often than other cells (38\% vs 28\%) \tabularnewline
\raggedright Track-stratified null & \raggedright 16/16 depleted and 16/16 enriched cells remain significant with unchanged direction \tabularnewline
\raggedright Two periods ($\leq$2023 vs $\geq$2024) & \raggedright 16/16 cells depleted in both periods; change does not differ from a permutation null ($p=.86$) \tabularnewline
\raggedright Keyword (stage-1) coding & \raggedright 13/16 depleted cells keep the same direction \tabularnewline
\raggedright Same-term self-pairing & \raggedright 33\% of records in the strongest enriched cell pair the same term; removing them drops 2 prediction-related depleted cells from the core \tabularnewline
\raggedright Rule-based keyword matrix & \raggedright 741 same-term records removed, 1,431 document edges; 13/16 same direction, 6 significant, 0 significantly reversed; only 7 cells reach the four-test core \tabularnewline
\bottomrule
\end{tabular}
\end{table}

\subsection{Study 2}

Table~\ref{tab:app_oeyear} gives the yearly intensities behind Figure~\ref{fig:patent_trend} and the synthetic calibration run before the real data. Table~\ref{tab:app_suite} summarises the calibration (P2), time-split (P3) and Gap Validity (P9) modules, and Table~\ref{tab:app_lift} gives the label face-validity lifts (P4). P4 is a title-keyword check, not a human coding audit.

\begin{table}[tbp]
\centering
\caption{OE(AI,DEC) by citing grant year, primary specification (applicant citations), with 95\% CI; and synthetic calibration of the aggregate rule (DIRECT; E3.1). The synthetic generator is simple (field-local citing with planted suppression), and its no-hole baseline lies far below the real OE.}
\label{tab:app_oeyear}
\footnotesize
\begin{tabular}{lrl@{\hspace{1.2cm}}lrl}
\toprule
Year & OE & 95\% CI & Year & OE & 95\% CI \\
\midrule
2013 & 3.90 & 3.16--4.85 & 2020 & 3.55 & 3.18--3.94 \\
2014 & 2.84 & 2.42--3.41 & 2021 & 3.34 & 3.03--3.69 \\
2015 & 3.79 & 3.14--4.55 & 2022 & 2.95 & 2.62--3.30 \\
2016 & 4.05 & 3.42--4.75 & 2023 & 3.05 & 2.78--3.34 \\
2017 & 3.29 & 2.78--3.86 & 2024 & 2.69 & 2.40--2.98 \\
2018 & 3.20 & 2.81--3.64 & 2025 & 2.34 & 2.08--2.64 \\
2019 & 3.99 & 3.48--4.54 & & & \\
\midrule
\multicolumn{6}{p{8.4cm}}{WLS slope $-0.097$ per year (95\% CI $-0.130$ to $-0.065$); last-3-year OE 2.69 vs reference median 1.095; label \texttt{PERSISTENT}.} \\
\midrule
\multicolumn{6}{l}{\emph{Synthetic calibration (before real data)}} \\
\multicolumn{6}{p{8.4cm}}{No hole: OE 0.388 vs reference lowest quartile 0.356; ratio 0.99 (0.79--1.21) $\rightarrow$ \texttt{HOLE\_\allowbreak NOT\_\allowbreak SUPPORTED}.} \\
\multicolumn{6}{p{8.4cm}}{30\% suppression: OE 0.275; ratio 0.66 (0.52--0.83) $\rightarrow$ detected.} \\
\multicolumn{6}{p{8.4cm}}{50\% suppression: OE 0.223; ratio 0.47 $\rightarrow$ detected.} \\
\bottomrule
\end{tabular}
\end{table}

\begin{table}[tbp]
\centering
\caption{Replication suite: calibration (P2), time split (P3) and Gap Validity (P9). Part~A is co-classification; Part~B is applicant citations (DIRECT; E3.3). The P2 null matrices are stub permutations that treat edges as independent, so they cannot detect anti-conservatism from document clustering. The first half of Part~A is small. The P9 entries are diagnostics, not evidence: the zero indirect term is an identity of bipartite graphs, and alternative routes are near-guaranteed in a dense $6\times7$ matrix.}
\label{tab:app_suite}
\footnotesize
\begin{tabular}{lp{6.0cm}ll}
\toprule
Module & Statistic & Part A & Part B \\
\midrule
P2 seeds & Seeds (20 $\times$ 10,000): cells always q $\leq$ .05 / boundary & 20 / 4 & 22 / 0 \\
P2 null & Mean false cells per matrix (1,000 null matrices) & 0.045 & 0.049 \\
 & Share of matrices with any false cell & 3.6\% & 4.4\% \\
P2 power & 20\% depletion detected, cells with $E\geq100$ & 18/27 = 67\% & 23/29 = 79\% \\
P2 segmentation & S1 inventive codes only: depleted (overlap with P1) & 8 (6) & 11 (8) \\
 & S2 single-code documents: depleted (overlap with P1) & 8 (5) & 8 (4) \\
\midrule
P3 time split & First half: documents, depleted & 2010--2017: 543, 0 & 2013--2018: 4,772, 3 \\
 & Second half: documents, depleted & 2018--2025: 4,773, 8 & 2019--2025: 26,071, 11 \\
 & Depleted cells replicated in both halves & 0 & 0 \\
\midrule
P9 Gap Validity & Indirect constraint term (bipartite; identity) & 0 & 0 \\
 & Alternative length-3 routes per depleted cell & 20--30 & 20--30 \\
 & Depleted cells at q $\leq$ .01 vs q $\leq$ .05 & 7 vs 9 & 9 vs 10 \\
\bottomrule
\end{tabular}
\end{table}

\begin{table}[tbp]
\centering
\caption{Label face validity (P4): lift of the title-keyword hit rate over other AI$\cup$DEC patents. A label passes if lift $\geq$ 1.5; all 11 pass (DIRECT; E3.3). This is not inter-rater reliability.}
\label{tab:app_lift}
\footnotesize
\begin{tabular}{lr@{\hspace{1.2cm}}lr}
\toprule
AI-method label & Lift & Decision-function label & Lift \\
\midrule
G06N3 & 4.69 & FORECAST\_OPTIMISE & 3.31 \\
G06N5 & 7.81 & RESOURCE\_SCHEDULING & 4.51 \\
G06N7 & 14.28 & WORKFLOW & 4.88 \\
G06N10 & 1,650 & RISK & 55.39 \\
G06N20 & 2.18 & STRATEGY (hit rate 4.5\%) & 4.40 \\
 & & PERFORMANCE & 6.32 \\
\bottomrule
\end{tabular}
\end{table}

\subsection{External-data feasibility}

Table~\ref{tab:app_sources} gives the scale and outcome of the six public sources screened for flow validity.

\begin{table}[tbp]
\centering
\caption{Six public sources screened for flow validity (DIRECT; E2). No source passes the independent-relation and validation gates.}
\label{tab:app_sources}
\footnotesize
\begin{tabular}{p{2.2cm}p{6.6cm}p{5.6cm}}
\toprule
Source & Scale & Relation type and outcome \tabularnewline
\midrule
\raggedright CORDIS H2020 & \raggedright 35,389 projects; 179,063 organisation participations; keyword candidates 195 projects $\times$ 1,105 organisations = 1,436 edges; projection 1,049 nodes, 12,181 edges, 102,660 within-project triads, 156 edges recurring across $\geq$2 projects & \raggedright Co-membership (affiliation). Degree-preserving bipartite swap null (49 reps): 5 testable candidates, $p=.41$--$.86$, BH q $\geq$ .80; 0 significant \tabularnewline
\raggedright GLEIF Level-2 & \raggedright 663,449 relationship records (481,247 active) & \raggedright Legal consolidation, not information flow \tabularnewline
\raggedright OpenAlex & \raggedright 8 works & \raggedright Fails gates \tabularnewline
\raggedright OCDS & \raggedright 0 verifiable buyer--supplier edges in the round-3 sample & \raggedright Fails gates \tabularnewline
\raggedright SEC & \raggedright Not fetched (requires a genuine contact User-Agent) & \raggedright Not assessed \tabularnewline
\raggedright PyPSA & \raggedright 8,786 buses & \raggedright Physical grid topology, not information flow \tabularnewline
\bottomrule
\end{tabular}
\end{table}

% ------------------------------------------------------------
\section{Multi-Level Programming Constraint Details}
\label{app:mlp}

All values in this appendix describe the simulated v1.0 prototype. They are design assumptions, not estimates, and describe no real firm. The model formalises only the organisational candidate explanation (Section~\ref{sec:discussion_hypothesis}), and its original anchor, the 16 V9 depleted cells, has been withdrawn.

\subsection{Full statement of the three levels}
\label{app:mlp_full}

This subsection states the model summarised in Section~\ref{sec:mlp_model}. Each process $i$ has an initial capability $X_i^0$ for every resource $X\in\mathcal{R}=\{D,A,R,G\}$.

\paragraph{Level 1: enterprise resource allocation (leader).}
\begin{align}
\max_{b}\quad & F_1 = 1.25\,\mathrm{RealizedValue} + 0.20\,\mathrm{AIAdoptions} - 0.22\,\mathrm{ResourceUse} - 0.75\,\mathrm{ResidualRisk} \label{eq:mlp_l1}\\
\text{s.t.}\quad & b_i^r\in\{0,1,2,3\}, \qquad \sum_{i\in I}\sum_{r\in\mathcal{R}} c_i^r\, b_i^r \leq B_{\mathrm{tot}}, \nonumber\\
& \bar{X}_i = \phi_X\!\left(X_i^0, b_i^X\right) = \min\!\left(3,\; X_i^0 + b_i^X\right), \qquad X\in\mathcal{R}, \nonumber\\
& (A_i,H_i,WI_i,EA_i)_{i\in I} \text{ solves Level 2 given } \bar{X}, \text{ and } q \text{ solves Level 3}. \nonumber
\end{align}
The toy map $\phi_X$ turns an allocation $b_i^X$ at unit cost $c_i^r$ into a capability ceiling $\bar{X}_i\in\{\bar D_i,\bar A_i,\bar R_i,\bar G_i\}$. RealizedValue, AIAdoptions and ResidualRisk are evaluated at the followers' responses.

\paragraph{Level 2: AI decision architecture and governance.} For each process the architecture chooses AI integration depth $A_i$, human oversight $H_i$, workflow integration $WI_i$ and assurance $EA_i$:
\begin{align}
\max_{(A_i,H_i,WI_i,EA_i)_{i\in I}}\quad & F_2 \label{eq:mlp_l2}\\
\text{s.t.}\quad & (A_i,H_i,WI_i,EA_i)\in\{0,\dots,3\}^4, \nonumber\\
& A_i\leq\bar{A}_i,\quad H_i\leq\bar{G}_i,\quad EA_i\leq\bar{G}_i,\quad WI_i\leq\bar{D}_i, \nonumber\\
& \bar{D}_i\geq D_{\mathrm{req}}(A_i),\quad \bar{R}_i\geq R_{\mathrm{req}}(A_i), \nonumber\\
& H_i\geq H_{\mathrm{req}}(A_i,\kappa_i),\quad EA_i\geq EA_{\mathrm{req}}(A_i,\kappa_i),\quad WI_i\geq WI_{\mathrm{req}}(A_i), \nonumber\\
& \sum_{i\in I}\left(0.42A_i + 0.24H_i + 0.24EA_i + 0.20WI_i\right)\leq K_\Gamma. \nonumber
\end{align}
The requirement thresholds are listed in Table~\ref{tab:app_thresholds}. The capacity constraint is shared by all processes; it represents the organisation's limited capacity to govern AI across processes. $F_2$ sums, over processes, weighted terms for feasibility, reliability and scalability (bounded linear indices of the capability surplus and the Level-2 state), minus integration cost and governance risk, plus the best Level-3 utility the architecture makes available. Its functional form and weights are specified in the canonical model file of the project repository and summarised below; the weights are not calibrated and carry no empirical content.

\paragraph{Level 3: human--AI configuration (followers).} Given its Level-2 state, each process chooses one configuration $m\in\mathcal{M}=\{$HumanOnly, AIAssist, HumanFirstAICheck, AIFirstReview, BoundedAutomation$\}$:
\begin{align}
\max_{q_{i\cdot}}\quad & \sum_{m\in\mathcal{M}} q_{im}\,U_{im}, \qquad U_{im} = \mathrm{Gain}_{im} - 0.45\,\mathrm{Burden}_{im} - 0.75\,\mathrm{ResidualRisk}_{im} \label{eq:mlp_l3}\\
\text{s.t.}\quad & q_{im}\in\{0,1\},\qquad \sum_{m\in\mathcal{M}} q_{im}=1,\qquad q_{im}\leq\mathrm{Allow}_{m\Gamma}, \nonumber\\
& q_{im}=1 \;\Rightarrow\; (A_i,H_i,WI_i,EA_i)\geq\rho_m \text{ componentwise}. \nonumber
\end{align}
The minimum requirement vectors $\rho_m=(A,H,WI,EA)$ are given in Table~\ref{tab:app_configs}; HumanOnly requires $(0,0,0,0)$, and the requirements rise to $(3,3,3,3)$ for BoundedAutomation.

\paragraph{Selection processes and the threshold $\tau$.} For a process that selects among alternatives, the AI component produces a score $s_i$, and $\tau_i$ is the score threshold above which it recommends the positive choice:
\begin{equation}
d_i=\mathbf{1}\!\left[s_i\geq\tau_i\right], \qquad v_i=\begin{cases}1 & \text{if } d_i \text{ matches the true outcome},\\ 0.35 & \text{otherwise},\end{cases}
\label{eq:mlp_tau}
\end{equation}
where $v_i$ scales the value the process realises. The threshold is a fixed design parameter of the process, not a decision variable of any level. In the prototype, $s_i$ and the true outcome are a single fixed realisation known to all levels, so $\tau_i$ is a sensitivity device, not a model of error rates.

\subsection{Governance regimes}

A governance regime $\Gamma$ is defined by the parameters in Table~\ref{tab:app_regimes} and by the permission matrix $\mathrm{Allow}_{m\Gamma}$. In the prototype source, the flexibility multiplier scales realised gain at Level~3, the risk multiplier scales residual risk at Level~3 and governance risk at Level~2, and the cost multiplier scales integration cost at Level~2. For general background on bilevel and multilevel programming, see \citet{dempe2002foundations}, \citet{colson2007overview}, \citet{fortunyamat1981representation}, \citet{lee2001fuzzy} and \citet{cheng2019fuzzy}; none of these sources supports a specific modelling choice of the prototype.

\begin{table}[tbp]
\centering
\caption{Governance regime parameters (SIMULATED; E4). The letters C, F and D are labels for parameter sets only and are not expanded. In the baseline, all five configurations are permitted under every regime.}
\label{tab:app_regimes}
\footnotesize
\begin{tabular}{lrrrr}
\toprule
Regime & Capacity $K_\Gamma$ & Risk multiplier & Cost multiplier & Flexibility multiplier \\
\midrule
C & 5.5 & 0.85 & 1.10 & 0.95 \\
F & 6.0 & 1.00 & 1.00 & 1.00 \\
D & 6.5 & 1.15 & 0.90 & 1.05 \\
\bottomrule
\end{tabular}
\end{table}

\subsection{Level-2 requirement thresholds and capacity}

Table~\ref{tab:app_thresholds} lists the threshold functions used in Eq.~\eqref{eq:mlp_l2}; $\kappa_i$ is the consequence class of process $i$. The shared capacity constraint is
\begin{equation}
\mathrm{CapUse}=\sum_{i\in I}\left(0.42A_i+0.24H_i+0.24EA_i+0.20WI_i\right)\leq K_\Gamma ,
\label{eq:app_capacity}
\end{equation}
with $\mathrm{CapUse}(0,0,0,0)=0$ and $K_\Gamma\geq0$. The prototype source also bounds $H_i$, $EA_i$ and $WI_i$ by $3A_i$. On the ordinal scale this restriction binds only at $A_i=0$, where it rules out oversight, assurance or workflow integration of an absent AI component. It is a design assumption of the prototype.

\begin{table}[tbp]
\centering
\caption{Level-2 requirement thresholds as functions of AI depth $A$ and consequence class $\kappa$ (canonical v1.0 specification; design assumptions, not calibrated).}
\label{tab:app_thresholds}
\footnotesize
\begin{tabular}{ll}
\toprule
Requirement & Definition \\
\midrule
Data, $D_{\mathrm{req}}(A)$ & $A$ \\
Reliability, $R_{\mathrm{req}}(A)$ & $A$ \\
Oversight, $H_{\mathrm{req}}(A,\kappa)$ & $0$ if $A=0$; otherwise $\min(3,\max(0,A+\kappa-2))$ \\
Assurance, $EA_{\mathrm{req}}(A,\kappa)$ & $0$ if $A=0$; otherwise $\min(3,\max(1,A+\kappa-2))$ \\
Workflow integration, $WI_{\mathrm{req}}(A)$ & $3$ for an execution process with $A=3$; otherwise $A$ \\
\bottomrule
\end{tabular}
\end{table}

\subsection{Structure of the Level-2 objective and the Level-3 terms}

The Level-2 objective $F_2$ and the value and risk terms are specified in full, with their weights, in the canonical model file of the project repository (\texttt{01\_\allowbreak CANONICAL\_\allowbreak MODEL\_\allowbreak v1\_\allowbreak 0.md}). We give their structure here so that the model can be read without that file; the weights are positive constants that have not been calibrated, and we do not reproduce them because they carry no empirical content. With uncalibrated weights $\omega_1,\dots,\omega_6>0$,
\begin{equation}
F_2=\sum_{i\in I}\Bigl[\omega_1\,\mathrm{Feas}_i+\omega_2\,\mathrm{Rel}_i+\omega_3\,\mathrm{Scal}_i-\omega_4\,\mathrm{IntCost}_i-\omega_5\,\mathrm{GovRisk}_i+\omega_6\max\!\bigl(0,U_i^{*}\bigr)\Bigr],
\label{eq:app_f2}
\end{equation}
where $U_i^{*}$ is the best Level-3 utility available to process $i$ in its Level-2 state. Feasibility $\mathrm{Feas}_i$, reliability $\mathrm{Rel}_i$ and scalability $\mathrm{Scal}_i$ are linear indices clipped to $[0,1]$. Feasibility increases in the capability surpluses $\bar D_i-D_{\mathrm{req}}$, $\bar R_i-R_{\mathrm{req}}$ and $\bar G_i-\max(H_{\mathrm{req}},EA_{\mathrm{req}})$. Reliability increases in $EA_i$, $H_i$ and $WI_i$ and decreases in $A_i\kappa_i$. Scalability increases in $\bar D_i$, $\bar R_i$ and $WI_i$. Integration cost is the regime's cost multiplier times a linear function increasing in $A_i$, $WI_i$, $H_i$ and $EA_i$. Governance risk is 0 when $A_i=0$; otherwise it is the regime's risk multiplier times a term proportional to $\kappa_i A_i$, divided by a term increasing in $H_i$ and $EA_i$.

At Level~3, $\mathrm{Gain}_{im}$ is the product of a configuration-specific base gain, a process value weight, a factor $\psi_i$ that increases in $EA_i$ and $WI_i$ and, for selection processes, is multiplied by the value factor $v_i$ of Eq.~\ref{eq:mlp_tau}, and the regime's flexibility multiplier. $\mathrm{Burden}_{im}$ is a configuration-specific base burden plus a linear term increasing in $H_i$, $EA_i$ and $WI_i$. $\mathrm{ResidualRisk}_{im}$ is $\kappa_i$ times a configuration-specific autonomy level times the regime's risk multiplier, divided by a term increasing in $H_i$ and $EA_i$. All configuration-specific constants are simulated.

\subsection{Level-3 configurations and the selection mechanism}

Table~\ref{tab:app_configs} gives the minimum Level-2 state each configuration requires. A process may choose configuration $m$ only if $\mathrm{Allow}_{m\Gamma}=1$ and its state meets $\rho_m$ componentwise; it then receives $U_{im}=\mathrm{Gain}_{im}-0.45\,\mathrm{Burden}_{im}-0.75\,\mathrm{ResidualRisk}_{im}$. For selection processes, the AI recommendation is $d_i=\mathbf{1}[s_i\geq\tau_i]$, and the realised value factor is 1 for a correct recommendation and 0.35 otherwise (Eq.~\ref{eq:mlp_tau}). Review configurations reduce risk or change utility but do not correct a wrong recommendation. Ties are broken deterministically in the order HumanOnly, AIAssist, HumanFirstAICheck, AIFirstReview, BoundedAutomation.

\begin{table}[tbp]
\centering
\caption{Minimum Level-2 requirements $\rho_m=(A,H,WI,EA)$ of the five human--AI configurations (canonical v1.0 specification).}
\label{tab:app_configs}
\footnotesize
\begin{tabular}{lcccc}
\toprule
Configuration $m$ & $A$ & $H$ & $WI$ & $EA$ \\
\midrule
HumanOnly & 0 & 0 & 0 & 0 \\
AIAssist & 1 & 0 & 1 & 1 \\
HumanFirstAICheck & 2 & 1 & 2 & 1 \\
AIFirstReview & 2 & 2 & 2 & 2 \\
BoundedAutomation & 3 & 3 & 3 & 3 \\
\bottomrule
\end{tabular}
\end{table}

\subsection{Existence and scope}

\paragraph{Lemma (existence of a selected tri-level profile).} Assume (i) all decision and state domains are finite; (ii) the zero allocation $b=0$ is admissible ($B_{\mathrm{tot}}\geq0$); (iii) HumanOnly has requirement $(0,0,0,0)$ and is permitted under every regime; (iv) Level-3 followers are separable conditional on upstream decisions; and (v) deterministic tie-breaking at every level. Then, for a fixed regime, initial state and parameter snapshot, at least one selected tri-level profile exists, and exhaustive backward induction terminates in finitely many steps.

\paragraph{Proof sketch.} \emph{Level~3:} HumanOnly is feasible for every Level-2 state, so each follower's feasible set is finite and non-empty; a maximiser exists, and tie-breaking selects one. \emph{Level~2:} for any capability vector the null architecture $A_i=H_i=WI_i=EA_i=0$ satisfies the ceilings, every threshold (all are 0 at $A=0$), the $A_i=0$ restriction and the capacity constraint, since $\mathrm{CapUse}=0\leq K_\Gamma$; the feasible set is finite and non-empty, so $F_2$ attains a maximum. \emph{Level~1:} $b=0$ is feasible and the allocation set is finite. Backward induction enumerates finite sets at every level and therefore terminates. $\square$

\paragraph{Scope.} The lemma establishes non-empty feasibility, existence of a selected response and finite termination for this bounded finite prototype only. A selected representative is not a unique optimum, and optimistic and pessimistic formulations are not generally equivalent \citep{dempe2002foundations}. Existence and exact enumeration do not imply tractability for larger, continuous or non-convex tri-level programmes \citep{jeroslow1985polynomial}. The Level-2 capacity constraint couples processes, but it is chosen by a single decision maker. If instead the Level-3 followers shared a constraint (for example, a common pool of reviewers), assumption (iv) would fail. A coupled or generalised equilibrium concept would then be required, and payoff interactions could rule out pure-strategy equilibria. Equilibrium existence in finite non-cooperative games \citep{nash1951noncooperative} serves only as a boundary reference; it does not cover coupled constraints.

\subsection{Conformance and mutation checks}

The solver was checked in two ways (E4). The \emph{conformance check} evaluates, for each of the three processes and every tuple of capability ceilings $(\bar D_i,\bar A_i,\bar R_i,\bar G_i)$, the canonical constraints written above independently of the solver's state generator and compares the two feasible sets. Across 768 cases it counted 3,121 canonical-feasible state occurrences and found 0 mismatches. The \emph{mutation check} runs the real generator with the $A_i=0$ restriction disabled and the corresponding branch relaxed. It detected a mismatch with the canonical set in 720/768 cases, which shows that the conformance check is sensitive to this class of error. Both checks share the threshold-function implementations, so neither can detect a common-mode defect in those functions. The total runtime of the solver script was 13.9 s. These checks verify the toy instance computationally; they leave every coefficient, threshold, score, outcome, permission matrix and regime ordering empirically unverified.

\subsection{Study 3: decision rules, stop rules and network test}
\label{app:study3}

This subsection states in full the rules summarised in Section~\ref{sec:mlp_study3}. No data have been collected, and all thresholds will be fixed in Study~3's own preregistration before data access.

\paragraph{Decision rules.} The organisational candidate is supported if, among episodes with a determinable binding constraint, internal permission, decision rights, governance capacity or allocation is the most frequent first binding constraint and exceeds the preregistered share, and if episodes that changed status after an internal rule or allocation changed did so without a recorded change in capability or data. Each alternative (data quality, cost, regulation and liability, task suitability) is supported by the same rule applied to its own signature in Table~\ref{tab:back_study3_rivals}. Internal decision rights and external regulation can both appear as a withheld permission; they are separated by whether the binding rule originates inside the firm, where management can change it, or outside it. If no candidate reaches the threshold, the result is reported as mixed, and the model is not calibrated as a governance-bottleneck model.

\paragraph{Stop rules.} The study stops and is reported as uninformative if the episodes with a determinable binding constraint fall below the preregistered minimum. Coding stops, and the codebook is revised as a dated amendment before any outcome analysis, if agreement between human coders on a pilot set falls below the preregistered minimum. If the organisational candidate is not supported, the model's governance layer is not calibrated, and the model is reported as unsupported.

\paragraph{Optional network test.} If de-identified metadata on tickets and meetings between AI staff and decision owners are available, Burt constraint \citep{burt1992structural} will be computed on real organisational ties, where embeddedness arguments locate economic action \citep{granovetter1985economic}, against a null that preserves both what each unit is responsible for and its degree.
